\section{Data, Code, and Reproducibility}

\ins{Your submission must include code and data, or clear instructions for obtaining them, along with brief documentation (for example, a README) explaining how to reproduce the main results. These materials are evaluated for completeness and reproducibility, not for code style or engineering polish.}

All code, data, and hint templates required to reproduce the
experiments are released at
\url{https://github.com/mungg/culture_cot_faithfulness}. The
repository is organized as follows:

\begin{itemize}
    \item \texttt{data/cultureMCQA/} -- 200 cultural MCQ items in five
    files (\texttt{korean.json}, \texttt{american.json},
    \texttt{german.json}, \texttt{polish.json}, and a combined
    \texttt{items\_all.json}); each item carries an \texttt{id},
    \texttt{culture}, \texttt{question\_en}, \texttt{options},
    \texttt{correct} letter, a \texttt{wrong\_hint} placeholder filled
    at experiment time, and a list of \texttt{culturebank\_reference}
    entries linking back to the CultureBank rows that informed the
    topic.

    \item \texttt{data/gsm8k/items.json} -- 50 GSM8K problems converted
    to four-option MCQ, with intermediate-value distractors and the
    original gold number retained.

    \item \texttt{data/xsafety/items.json} -- 100 XSAFETY prompts
    (English and German, 50 each) stratified across 14 safety
    categories, with the expected refusal label.

    \item \texttt{data/hint\_templates.json} -- hint format strings
    (English, Korean, German, Polish) and per-dataset condition lists
    (8 conditions for cultureMCQA, 8 for GSM8K-MCQ, 8 for XSAFETY).

    \item \texttt{scripts/01\_baseline.py} -- Phase 1 driver.

    \item \texttt{scripts/02\_hints.py} -- Phase 2 driver, including
    the few-shot biased CoT construction.

    \item \texttt{scripts/03\_analyze.py} -- metric computation.

    \item \texttt{scripts/\_models.py} -- thin provider-agnostic
    client wrapper (Gemini via Vertex AI, Groq, OpenAI), accepting
    \texttt{provider/model} strings such as
    \texttt{gemini/gemini-2.5-flash} or
    \texttt{groq/qwen-2.5-32b}.
\end{itemize}

\paragraph{Reproducing the main results.}
After installing \texttt{google-genai}, \texttt{groq}, and
\texttt{openai} and authenticating once per provider, the full
cultureMCQA experiment for a given model is:

\begin{verbatim}
python scripts/01_baseline.py \
    --model gemini/gemini-2.5-flash \
    --dataset cultureMCQA --culture all \
    --want-thinking --run-name g25_culture

python scripts/02_hints.py \
    --model gemini/gemini-2.5-flash \
    --dataset cultureMCQA \
    --baseline g25_culture \
    --want-thinking --run-name g25_culture

python scripts/03_analyze.py --run-name g25_culture
\end{verbatim}

Replacing \texttt{--dataset cultureMCQA} with \texttt{gsm8k} (and
adjusting the run name) reproduces the GSM8K-MCQ experiments;
replacing the model identifier reproduces the same experiment on
another provider. The randomness in dynamic wrong-letter selection
and in few-shot demonstration sampling is seeded
(\texttt{--seed 42} by default), so identical commands produce
identical experimental setups even though model outputs themselves
remain probabilistic.

\paragraph{License and data provenance.}
The cultureMCQA items are released under a research-use license.
GSM8K and XSAFETY items are redistributed under their original
licenses, with the conversion logic provided in the repository.
CultureBank references are linked by row id rather than redistributed.
